% Worked companion to G11_T1_L2_C8C9C10_exam.tex.
% Content: source exam. Presentation: template_exam_solution.tex.
\makeatletter
\def\input@path{{./}{../}{../../}{../../../}{../../../../}{preamble/}{../preamble/}{../../preamble/}{../../../preamble/}{../../../../preamble/}}
\makeatother
\input{preamble_exams.tex}
\graphicspath{{./}{../}{../../}{../../../}{../../../../}{preamble/}{../preamble/}{../../preamble/}{../../../preamble/}{../../../../preamble/}}
\linespread{1.0}
\setlength{\parskip}{5pt}
\usepackage{needspace}
\title{Grade 11 --- Learning Evidence 2\\C8, C9 and C10 --- Exam Solutions\\\large 1st Term \quad 2026--2027}
\author{}
\date{}
\begin{document}
\maketitle
\setcounter{tocdepth}{1}
\setcounter{secnumdepth}{2}
\phantomsection
\label{toc}
\tableofcontents

\section*{Learning objective}
Apply continuous probability models and sampling distributions to investigate economic and financial questions, justify statistical methods, and communicate evidence-based conclusions.

\section*{Criteria assessment}
Show all relevant calculations and clearly justify your answers. Partial credit may be awarded when the correct procedure is clearly demonstrated, even if arithmetic errors are made. No points will be awarded for answers without a demonstrated procedure, except for the brief research selection in Question 7.

A minimum of 4 points is required to pass each criterion. A 1-point penalty will be applied to each criterion if the student's name is not written on the assessment or if the assessment is submitted late.

\begin{center}
\begin{tabular}{|c|c|c|c|}
\hline
\textbf{Question} & \textbf{C8} & \textbf{C9} & \textbf{C10} \\ \hline
1 & 2 & -- & -- \\ \hline
2 & 2 & -- & -- \\ \hline
3 & 2 & -- & -- \\ \hline
4 & -- & 2 & -- \\ \hline
5 & -- & 2 & -- \\ \hline
6 & -- & 2 & -- \\ \hline
7 & -- & -- & 6 \\ \hline
\textbf{Total} & \textbf{6} & \textbf{6} & \textbf{6} \\ \hline
\end{tabular}
\end{center}

Answer on separate paper and show your reasoning. Use the supplied normal and Student's t reference sheets.
Report Student's t critical values and normal critical values from the infinity row to \textbf{3 decimal places}, and normal probabilities to \textbf{4 decimal places}.
For Student's t cutoffs between printed entries, give probability bounds without interpolation.
Each criterion is worth \textbf{6 points} (\textbf{18 points} total); at least \textbf{4 points per criterion} are required to pass that criterion.

Questions 1--6 concern independent random observations from normal populations.
All critical values and probability questions refer to the sample mean standardized relative to the population mean; critical values have no units.



\section*{Evaluated criteria}
\begin{enumerate}[label=\textbf{C\arabic*},start=8,leftmargin=3em]
\item Uses Student's t-distribution tables, degrees of freedom, and one-tailed or two-tailed areas to determine critical values, central probabilities, and tail probabilities or probability ranges.
\item Calculates and compares critical values and probability areas from the standard normal and Student's t-distributions, selecting and justifying the appropriate distribution according to whether the population standard deviation is known or estimated from the sample.
\item Integrates continuous probability modelling, simulation, sampling analysis, and portfolio-risk measures to investigate an economic or financial question, justifying the selected methods and communicating results, assumptions, limitations, and evidence-based conclusions.
\end{enumerate}

\section*{Models and table conventions}
For Questions 1--6, the observations are independent and come from normal populations.
The sample mean $\bar X$ is therefore normal. If the population standard deviation $\sigma$ is known,
\[
 Z=\frac{\bar X-\mu}{\sigma/\sqrt n}\sim N(0,1).
\]
If $\sigma$ is unknown and estimated by the sample standard deviation $S$, then
\[
 T=\frac{\bar X-\mu}{S/\sqrt n}\sim t_\nu,\qquad \nu=n-1.
\]
Here $S$ varies from sample to sample; a reported sample value is an estimate of $\sigma$.
These models are exact under the stated normal-population assumptions, including for small samples.
All critical values below are dimensionless.

The year-local reference sheets are in \texttt{book2026-2027/fundamentals/normal-table.tex}
(and its PDF), also linked by the C8 and C9 practice activities.
The normal sheet gives cumulative areas $\Phi(z)=P(Z\le z)$, not areas from $0$ to $z$.
Its positive-$z$ entries have four decimal places.
The Student's $t$ sheet gives three-decimal positive critical values with the convention
\[
 P(T_\nu>t_{\alpha,\nu})=\alpha,\qquad
 P(T_\nu<-t_{\alpha,\nu})=\alpha,\qquad
 P(-t_{\alpha,\nu}<T_\nu<t_{\alpha,\nu})=1-2\alpha.
\]
Its printed columns are \emph{one-tail} probabilities
$0.30,0.20,0.10,0.05,0.025,0.01,0.005$.
A central area $c$ corresponds to combined two-tail area $1-c$ and one-tail area $(1-c)/2$;
there is no separate two-tail column on this sheet.
The $\infty$ row gives the normal critical values.
Table critical values are rounded; probability equalities involving those printed cutoffs
are understood at table precision. Bounds between printed $t$ entries use no interpolation.

\needspace{12\baselineskip}
\sectionScored{Export inspection costs}{C8-\sml{2}}{sec:q1}
{
An exporter studies average inspection costs using 12 shipments. The population standard deviation is unknown and estimated from the sample.
\QuestionsTasks{
{8}/{sol:q1a}/{\textbf{(a)} Determine the degrees of freedom and the positive critical value that leaves 1 percent of the probability to its right. \score{1}},
{8}/{sol:q1b}/{\textbf{(b)} Determine the negative critical value below which 5 percent of the probability lies. Justify its sign and the table column used. \score{1}}
}

\needspace{6\baselineskip}
\phantomsection
\label{sol:q1a}
\subsection*{C8 -- (a) Degrees of freedom and right-tail cutoff \score{1}}

The population standard deviation is unknown. With $n=12$ independent observations from a normal population,
$T=(\bar X-\mu)/(S/\sqrt{12})\sim t_{11}$, because $\nu=12-1=11$.
We want $P(T_{11}>c)=0.01$.
In row $\nu=11$, the one-tail $\alpha=0.01$ column (equivalent to two tails totalling $0.02$) gives
\[
 \boxed{\nu=11,\qquad c=t_{0.01,11}=2.718.}
\]
The cutoff is positive because only the upper $1\%$ lies beyond it.

\secInicio[sec:q1]

\needspace{6\baselineskip}
\phantomsection
\label{sol:q1b}
\subsection*{C8 -- (b) Left-tail cutoff \score{1}}

We want $P(T_{11}<c)=0.05$. Row $11$, one-tail column $\alpha=0.05$, gives $t_{0.05,11}=1.796$.
By symmetry about zero,
\[
 P(T_{11}<-t_{0.05,11})=P(T_{11}>t_{0.05,11})=0.05,
 \qquad \boxed{c=-1.796.}
\]
The table prints the positive upper-tail cutoff; its negative is the lower-tail cutoff.
The corresponding combined two-tail area is $0.10$.

\secInicio[sec:q1]
}

\needspace{12\baselineskip}
\sectionScored{Market-stall electricity use}{C8-\sml{2}}{sec:q2}
{
A study investigates average electricity use across 25 market stalls. The population standard deviation is unknown and estimated from the sample.
\QuestionsTasks{
{8}/{sol:q2a}/{\textbf{(a)} Determine the degrees of freedom. For a central region containing 98 percent of the probability, determine the combined probability outside it and the probability in each tail. \score{1}},
{8}/{sol:q2b}/{\textbf{(b)} Determine the 2 critical values, symmetric about 0, enclosing this central region. Identify the table column used. \score{1}}
}

\needspace{6\baselineskip}
\phantomsection
\label{sol:q2a}
\subsection*{C8 -- (a) Degrees of freedom and tail areas \score{1}}

With unknown $\sigma$ estimated by $S$ and $n=25$ independent normal observations,
$T=(\bar X-\mu)/(S/\sqrt{25})\sim t_{24}$.
\[
 \boxed{\nu=25-1=24,\qquad P(\text{outside})=1-0.98=0.02,
 \qquad P(\text{each tail})=0.02/2=0.01.}
\]
Symmetry splits the combined outside area equally between the two tails.

\secInicio[sec:q2]

\needspace{6\baselineskip}
\phantomsection
\label{sol:q2b}
\subsection*{C8 -- (b) Central 98 percent cutoffs \score{1}}

We seek $P(-c<T_{24}<c)=0.98$, so $P(T_{24}>c)=0.01$.
Row $24$, one-tail column $\alpha=0.01$ (combined two-tail area $0.02$), gives $c=2.492$.
\[
 \boxed{\text{Critical values: }-2.492\text{ and }2.492.}
\]
The table-based central area is $1-2(0.01)=0.98$.

\secInicio[sec:q2]
}

\needspace{12\baselineskip}
\sectionScored{Cold-storage energy costs}{C8-\sml{2}}{sec:q3}
{
An analyst studies average energy costs using 15 cold-storage facilities. The population standard deviation is unknown and estimated from the sample.
\QuestionsTasks{
{8}/{sol:q3a}/{\textbf{(a)} Determine the degrees of freedom and give bounds for the probability that the standardized sample mean exceeds positive 2.40. \score{1}},
{8}/{sol:q3b}/{\textbf{(b)} Give bounds for the probability that the standardized sample mean is below negative 2.40, and for the probability that it lies between negative 2.40 and positive 2.40. \score{1}}
}

\needspace{6\baselineskip}
\phantomsection
\label{sol:q3a}
\subsection*{C8 -- (a) Degrees of freedom and right-tail bounds \score{1}}

Unknown $\sigma$ is estimated from $n=15$ independent normal observations, so
$T=(\bar X-\mu)/(S/\sqrt{15})\sim t_{14}$, with $\nu=15-1=14$.
The adjacent entries in row $14$ are
\[
 t_{0.025,14}=2.145<2.40<2.624=t_{0.01,14}.
\]
The right-tail area decreases as the positive cutoff increases. Hence
\[
 P(T_{14}>2.624)<P(T_{14}>2.40)<P(T_{14}>2.145),
\]
\[
 \boxed{\nu=14,\qquad 0.01<P(T_{14}>2.40)<0.025.}
\]
These are the tightest bounds from the adjacent columns available on the supplied sheet.
No interpolation is used.

\secInicio[sec:q3]

\needspace{6\baselineskip}
\phantomsection
\label{sol:q3b}
\subsection*{C8 -- (b) Left-tail and central bounds \score{1}}

Write $p=P(T_{14}>2.40)$, where $0.01<p<0.025$.
By symmetry $P(T_{14}<-2.40)=p$, so
\[
 \boxed{0.01<P(T_{14}<-2.40)<0.025.}
\]
The two excluded tails have combined probability $2p$, giving
\[
 0.02<2p<0.05.
\]
Multiplication by $-1$ reverses the inequalities; adding $1$ gives
\[
 -0.02>-2p>-0.05
 \quad\Longrightarrow\quad
 0.98>1-2p>0.95.
\]
Thus, in increasing order,
\[
 \boxed{0.95<P(-2.40<T_{14}<2.40)<0.98.}
\]
These are table bounds, not an exact numerical $t$ probability.

\secInicio[sec:q3]
}

\needspace{12\baselineskip}
\sectionScored{Bottling-line measurements}{C9-\sml{2}}{sec:q4}
{
A bottling plant studies average bottle volume using 16 bottles. The population standard deviation is known to be 3 millilitres.
\QuestionsTasks{
{9}/{sol:q4a}/{\textbf{(a)} Choose and justify the distribution model for the standardized sample mean. Determine the negative critical value below which 2.5 percent of the probability lies. \score{1}},
{9}/{sol:q4b}/{\textbf{(b)} Calculate the probability that the standardized sample mean lies between negative 1.00 and positive 1.50. \score{1}}
}

\needspace{6\baselineskip}
\phantomsection
\label{sol:q4a}
\subsection*{C9 -- (a) Model and lower critical value \score{1}}

The population is normal and $\sigma=3$ millilitres is known. Therefore
\[
 Z=\frac{\bar X-\mu}{3/\sqrt{16}}=\frac{\bar X-\mu}{0.75}\sim N(0,1).
\]
The standard error is $0.75$ millilitres; $Z$ has no units.
Normality makes this model exact even though $n=16$ is small.
We need $P(Z<c)=0.025$. The $\infty$ row, one-tail column $\alpha=0.025$
(equivalent to combined two-tail area $0.05$), gives $1.960$.
By symmetry,
\[
 P(Z<-1.960)=P(Z>1.960)\approx0.0250,
 \qquad \boxed{c=-1.960.}
\]

\secInicio[sec:q4]

\needspace{6\baselineskip}
\phantomsection
\label{sol:q4b}
\subsection*{C9 -- (b) Probability between the cutoffs \score{1}}

The normal sheet gives $\Phi(1.50)=0.9332$ (row $1.5$, column $0.00$)
and $\Phi(1.00)=0.8413$ (row $1.0$, column $0.00$).
Its entries are already cumulative, so no addition of $0.5$ is needed.
Symmetry gives
\[
 \Phi(-1.00)=1-\Phi(1.00)=1-0.8413=0.1587.
\]
Subtract the cumulative area below the lower endpoint:
\[
 P(-1.00<Z<1.50)=\Phi(1.50)-\Phi(-1.00)
 \approx0.9332-0.1587=\boxed{0.7745}.
\]
If using a different table of areas from $0$ to $z$, the equivalent calculation would be
$(0.8413-0.5000)+(0.9332-0.5000)=0.7745$.

\secInicio[sec:q4]
}

\needspace{12\baselineskip}
\sectionScored{Repair-time averages}{C9-\sml{2}}{sec:q5}
{
A workshop studies average repair time using 10 repairs. The population standard deviation is unknown; the sample standard deviation is 4 minutes.
\QuestionsTasks{
{9}/{sol:q5a}/{\textbf{(a)} Choose and justify the distribution model, determine the degrees of freedom, and find the positive critical value that the standardized sample mean exceeds with probability 2.5 percent. \score{1}},
{9}/{sol:q5b}/{\textbf{(b)} Estimate the probability that the standardized sample mean lies between negative 1.833 and positive 2.821. \score{1}}
}

\needspace{6\baselineskip}
\phantomsection
\label{sol:q5a}
\subsection*{C9 -- (a) Model and upper critical value \score{1}}

The population standard deviation is unknown; $4$ minutes is the observed sample standard deviation,
not a known $\sigma$. For $n=10$ independent normal observations,
\[
 T=\frac{\bar X-\mu}{S/\sqrt{10}}\sim t_9,\qquad \nu=10-1=9.
\]
We seek $P(T_9>c)=0.025$. Row $9$, one-tail column $\alpha=0.025$
(equivalent to combined two-tail area $0.05$), gives
\[
 \boxed{\nu=9,\qquad c=t_{0.025,9}=2.262.}
\]

\secInicio[sec:q5]

\needspace{6\baselineskip}
\phantomsection
\label{sol:q5b}
\subsection*{C9 -- (b) Table-based interval estimate \score{1}}

In row $9$, the one-tail columns $0.05$ and $0.01$ print
\[
 t_{0.05,9}=1.833,\qquad t_{0.01,9}=2.821.
\]
At table precision, symmetry gives $P(T_9<-1.833)\approx0.05$,
and the upper endpoint gives $P(T_9>2.821)\approx0.01$.
Subtract both disjoint tails:
\[
 P(-1.833<T_9<2.821)
 \approx1-0.05-0.01=\boxed{0.94\text{ (table-based estimate)}}.
\]
The endpoints are rounded printed critical values. Thus $0.94$ is the intended table-based estimate,
not an exact probability for those rounded endpoints. No interpolation is needed.

\secInicio[sec:q5]
}

\needspace{12\baselineskip}
\sectionScored{Comparing parcel-mass studies}{C9-\sml{2}}{sec:q6}
{
2 studies investigate average parcel mass. Study A uses 12 parcels and a known population standard deviation of 0.4 kilograms. Study B uses 12 parcels; its population standard deviation is unknown and estimated by a sample standard deviation of 0.4 kilograms.
\QuestionsTasks{
{9}/{sol:q6a}/{\textbf{(a)} Justify the distribution model for each study's standardized sample mean. Determine each model's 2 critical values, symmetric about 0, enclosing the central 95 percent of the probability. Calculate the difference between the positive critical values and explain it in terms of tail behaviour. \score{1}},
{9}/{sol:q6b}/{\textbf{(b)} Study B increases its sample to 61 parcels, with population standard deviation still unknown. State its model and degrees of freedom, and determine its new critical values enclosing the central 95 percent of the probability. Explain the trend toward the normal model and why sample size alone does not determine the model choice. \score{1}}
}

\needspace{6\baselineskip}
\phantomsection
\label{sol:q6a}
\subsection*{C9 -- (a) Comparing the two models \score{1}}

For $12$ independent normal observations, Study A knows $\sigma=0.4$ kilograms;
Study B uses $S$, observed as $0.4$ kilograms. Thus
\[
 Z_A=\frac{\bar X_A-\mu_A}{0.4/\sqrt{12}}\sim N(0,1),
 \qquad T_B=\frac{\bar X_B-\mu_B}{S/\sqrt{12}}\sim t_{11}.
\]
The same numerical standard deviation does not mean the same information:
$\sigma$ is a population constant; $S$ varies between samples.
For a central $95\%$, the outside area is $0.05$, or $0.025$ per tail.
Use the one-tail $\alpha=0.025$ column:
\begin{center}
\begin{tabular}{llll}
\toprule
Study & Model / row & Lower cutoff & Upper cutoff\\
\midrule
A & $N(0,1)$ / $\infty$ & $-1.960$ & $1.960$\\
B & $t_{11}$ / $11$ & $-2.201$ & $2.201$\\
\bottomrule
\end{tabular}
\end{center}
Each central area is $1-2(0.025)=0.95$ by symmetry and complements.
Using the displayed three-decimal cutoffs, the difference between positive cutoffs is
\[
 \boxed{2.201-1.960=0.241.}
\]
Student's $t$ has heavier tails because estimating $\sigma$ adds uncertainty.
Its cutoff must be farther from zero to leave only $2.5\%$ in the upper tail.

\secInicio[sec:q6]

\needspace{6\baselineskip}
\phantomsection
\label{sol:q6b}
\subsection*{C9 -- (b) Increasing Study B to 61 parcels \score{1}}

With unknown $\sigma$ and $61$ independent normal observations,
\[
 T_B=\frac{\bar X_B-\mu_B}{S/\sqrt{61}}\sim t_{60},
 \qquad \boxed{\nu=61-1=60.}
\]
Row $60$, one-tail column $\alpha=0.025$, gives $2.000$. By symmetry, the central $95\%$ cutoffs are
\[
 \boxed{\text{Critical values: }-2.000\text{ and }2.000.}
\]
The central area is $1-2(0.025)=0.95$ at table precision. The positive cutoff moves from $2.201$ toward $1.960$;
the normal-cutoff gap is now $2.000-1.960=0.040$. With more degrees of freedom, $S$ estimates $\sigma$ more reliably
and the $t$ tails approach the normal tails. Sample size does not make $\sigma$ known: the exact model remains $t_{60}$.

\secInicio[sec:q6]
}

\needspace{12\baselineskip}
\sectionScored{Choose your research option}{C10-\sml{6}}{sec:q7}
{
\textbf{During the exam:} select \textbf{1} of the following options and write a very brief description (1--2 sentences) of real-life data that could plausibly follow your selected distribution. \score{6}
\begin{enumerate}[label=\textbf{(\alph*)},leftmargin=*]
\item A continuous variable modelled by a linear probability density over a stated range.
\item Waiting times or intervals between events modelled by an exponential distribution.
\item Normally distributed measurements.
\end{enumerate}
No datasets, sources, calculations, parameter estimates, or detailed model justification are required during the exam.

\textbf{Online report submission:} your PDF report will focus on the \textbf{1} selected option and must be submitted virtually through \textbf{Open}. This online submission is separate from the brief response written during the exam. To demonstrate the full C10 criterion, investigate an economic or financial question using real data and your selected continuous model. Include simulation, sampling analysis, and portfolio-risk measures; justify the selected methods and communicate results, assumptions, limitations, and evidence-based conclusions. The 6 points shown here assess the brief exam response; the report provides the broader C10 evidence.
\QuestionsTasks{
{10}/{sol:q7a}/{Select only one option and give a brief description (1--2 sentences) of plausible real-life data. \score{6}}
}

\needspace{6\baselineskip}
\phantomsection
\label{sol:q7a}
\subsection*{C10 -- Acceptable alternative brief responses \score{6}}

This selection task has no unique answer. A student submits \textbf{only one} 1--2 sentence response,
choosing one option. The following are \textbf{alternatives}, not three required submissions:
\begin{enumerate}[label=\textbf{(\alph*)},leftmargin=*]
\item \textbf{Linear density.} I would study customers' spending in thousands of Colombian pesos on a promotional service priced continuously between 20 and 60 thousand pesos, proposing a linearly decreasing density over $[20,60]$. I would check whether real observations support that proposed model.
\item \textbf{Exponential waiting times.} I would study the minutes between payment requests reaching a small business during a stable trading period, proposing an exponential waiting-time model if arrivals are independent and occur at a constant average rate. I would check that assumption with real observations.
\item \textbf{Normal measurements.} I would study the fill volumes in millilitres of bottles sold by a beverage producer, proposing a normal model during stable production. I would check whether real measurements support that model.
\end{enumerate}
Each describes possible real-life economic data and a \emph{proposed} model requiring validation;
none claims that uncollected data follow the model. Other plausible responses are acceptable.
The brief exam response requires no dataset, sources, calculations, parameter estimates, or detailed model justification.
The original six points apply to this brief selection; no additional point allocation is imposed here.

\textbf{Separate online PDF report (not the brief exam response).}
Submit through \textbf{Open}, focusing on the \textbf{one selected option}.
Investigate an economic or financial question using real data and the selected continuous model.
Include simulation, sampling analysis, and portfolio-risk measures; justify methods and present results,
assumptions, limitations, and evidence-based conclusions. These are the broader C10 requirements;
no completed report, dataset, simulation, portfolio analysis, or empirical conclusion is supplied by these examples.

\secInicio[sec:q7]
}

\end{document}
